Naturally, a constraint solver uses a BTS algorithm within propagation. At each search decision, so at every assignment, it propagates the whole set of
constraints to remove inconsistent values from the domains of the future variables. \vspace{7pt}

The main advantages of a mechanism like this are:
\begin{enumerate}
    \item \textbf{Shrink} the domains of the future variables.
    \item \textbf{Propagate} all the constraints before search, and only when necessary ones during the search.
\end{enumerate}

However, both approaches, respectively with and without propagation, have the same complexity, equal to: 
$$ O(d^n) $$
where:
\begin{itemize}
    \renewcommand{\labelitemi}{-}
    \item $d \rightarrow$ domain size
    \item $n \rightarrow$ number of decision variables
\end{itemize}
For sure, using the propagation of constraints we can reduce the complexity of BTS algorithm, even though it remains exponential in size.
\begin{example}
    e.g. Same BTS example as before but with propagation. \vspace{7pt}

    Before the search, or the assignments, we start from the propagation, observing which values have a support in the other future variables domains.
    $$ C_2 : D(X_2) = \{1, 2, \bar{3}\}, D(X_3) = \{\bar{1}, 2, 3\}$$
    After the initial propagation, we can already state that values $3$ from $D(X_2)$ and $1$ from $D(X_3)$ do not have a support. Therefore, we can erase them 
    before the search phase. \vspace{7pt}

    \textbf{Why don't we have also constraint $C_1$}? 

    Because $C_1$ is already in AC, so we don't need to propagate. However, after first assignment ($X_1 = 1$) something changes. Since we assign a value 
    to the decision variable $X_1$, we have to propagate the first constraint and the second one since they share a specific decision variable $X_2$. 
    $$ C_1 : D(X_2) = \{\bar{1}, 2\}, C_2 : D(X_3) = \{\bar{2}, 3\} $$
\end{example}